
\documentclass[12pt,a4paper]{article}

\usepackage{cmap}
\usepackage[T2A,T1]{fontenc}
\usepackage[utf8]{inputenc}
\usepackage[russian]{babel}
\usepackage{tempora}
\usepackage{geometry}
\usepackage{amsmath,amssymb}
\usepackage{array}
\usepackage{booktabs}
\usepackage{longtable}
\usepackage{xcolor}
\usepackage{listings}
\usepackage{verbatim}
\usepackage{caption}
\usepackage{setspace}
\usepackage{hyperref}

\geometry{left=25mm,right=15mm,top=20mm,bottom=20mm}
\onehalfspacing
\hypersetup{colorlinks=true,linkcolor=black,urlcolor=black}

\lstset{
    language=Python,
    basicstyle=\fontencoding{T1}\ttfamily\scriptsize,
    numbers=left,
    numberstyle=\fontencoding{T1}\tiny,
    stepnumber=1,
    numbersep=7pt,
    frame=single,
    breaklines=true,
    breakatwhitespace=false,
    columns=fullflexible,
    keepspaces=true,
    showstringspaces=false,
    tabsize=4,
    keywordstyle=\bfseries,
    commentstyle=\itshape,
    stringstyle=\ttfamily,
    captionpos=b
}

\begin{document}

\begin{titlepage}
    \centering
    {\large Методы оптимизации\par}
    \vspace{1.2cm}
    {\Large\bfseries Лабораторная работа №6\par}
    \vspace{0.4cm}
    {\large Тема: «Методы поиска глобального экстремума»\par}
    \vspace{0.8cm}
    {\large Вариант 11\par}
    \vfill
    \begin{flushright}
        Выполнили: Григорьев Д.А., Столбченко И.В.
    \end{flushright}
    \vfill
    {Санкт-Петербург\par}
    {2026\par}
\end{titlepage}

\tableofcontents
\newpage

\section{Цель работы}
Цель лабораторной работы --- решить две оптимизационные задачи с помощью эвристических методов:
\begin{enumerate}
    \item задачу коммивояжера с помощью генетического алгоритма;
    \item задачу поиска кратчайшего пути от вершины A до вершины G с помощью алгоритма муравьиной колонии.
\end{enumerate}

\section{Задание 1. Задача коммивояжера}
\subsection{Постановка задачи}
Дано 5 городов и матрица расстояний между ними. Требуется найти замкнутый маршрут минимальной длины, который проходит через каждый город ровно один раз и возвращается в начальный город. В работе начальным городом считается город 1.

В задании требуется использовать генетический алгоритм. Размер популяции задан как $N=4$, вероятность мутации равна $0{,}01$. Мутация реализуется как случайная перестановка двух генов в геноме.

\subsection{Исходные данные}
Для варианта 11 используется следующая матрица расстояний:

\begin{table}[h!]
\centering
\caption{Матрица расстояний для задачи коммивояжера, вариант 11}
\begin{tabular}{c|ccccc}
\toprule
 & Город 1 & Город 2 & Город 3 & Город 4 & Город 5 \\
\midrule
Город 1 & 0  & 1  & 11 & 10 & 5 \\
Город 2 & 1  & 0  & 11 & 6  & 10 \\
Город 3 & 11 & 11 & 0  & 6  & 1 \\
Город 4 & 10 & 6  & 6  & 0  & 1 \\
Город 5 & 5  & 10 & 1  & 1  & 0 \\
\bottomrule
\end{tabular}
\end{table}

\subsection{Суть генетического алгоритма}
Генетический алгоритм имитирует естественный отбор. Каждая особь популяции представляет один возможный маршрут. В данной работе геномом является перестановка городов со 2-го по 5-й, так как город 1 зафиксирован как начальный и конечный.

Целевая функция --- длина маршрута:
\[
L = \sum_{i=1}^{n} d(v_i, v_{i+1}), \quad v_{n+1}=v_1.
\]
Так как нужно минимизировать длину, функция приспособленности взята в виде:
\[
F = \frac{1}{L}.
\]

Основные шаги алгоритма:
\begin{enumerate}
    \item создается начальная популяция из 4 случайных маршрутов;
    \item для каждой особи вычисляется длина маршрута и приспособленность;
    \item родители выбираются методом рулетки: чем меньше длина маршрута, тем выше шанс выбора;
    \item применяется упорядоченное скрещивание, чтобы получить корректную перестановку городов;
    \item с вероятностью $0{,}01$ выполняется мутация --- обмен двух случайных городов;
    \item лучшая особь сохраняется без изменений за счет элитизма.
\end{enumerate}

\section{Задание 2. Алгоритм муравьиной колонии}
\subsection{Постановка задачи}
Требуется найти кратчайший путь из вершины A в вершину G. Граф задан матрицей весов. Значение $\infty$ означает, что ребра между соответствующими вершинами нет.

\subsection{Исходные данные}
\begin{table}[h!]
\centering
\caption{Матрица весов графа для алгоритма муравьиной колонии, вариант 11}
\begin{tabular}{c|ccccccc}
\toprule
 & A & B & C & D & E & F & G \\
\midrule
A & $\infty$ & 17 & 5 & $\infty$ & 8 & $\infty$ & $\infty$ \\
B & 17 & $\infty$ & $\infty$ & $\infty$ & $\infty$ & $\infty$ & $\infty$ \\
C & 5 & $\infty$ & $\infty$ & 17 & 17 & $\infty$ & 45 \\
D & $\infty$ & $\infty$ & 17 & $\infty$ & 9 & 11 & 22 \\
E & 8 & $\infty$ & 17 & 9 & $\infty$ & $\infty$ & $\infty$ \\
F & $\infty$ & $\infty$ & $\infty$ & 11 & $\infty$ & $\infty$ & 13 \\
G & $\infty$ & $\infty$ & 45 & 22 & $\infty$ & 13 & $\infty$ \\
\bottomrule
\end{tabular}
\end{table}

\subsection{Суть алгоритма муравьиной колонии}
Алгоритм муравьиной колонии моделирует поведение муравьев, которые выбирают путь с учетом длины ребра и количества феромона на нем. Чем короче ребро и чем больше на нем феромона, тем выше вероятность его выбора.

Вероятность перехода из вершины $i$ в вершину $j$ вычисляется по формуле:
\[
P_{ij}=\frac{\tau_{ij}^{\alpha}\eta_{ij}^{\beta}}{\sum\limits_{k \in S}\tau_{ik}^{\alpha}\eta_{ik}^{\beta}},
\]
где $\tau_{ij}$ --- количество феромона на ребре, $\eta_{ij}=1/d_{ij}$ --- эвристическая привлекательность ребра, $S$ --- множество доступных соседних вершин, $\alpha$ и $\beta$ --- параметры влияния феромона и расстояния.

После прохождения муравьев феромон испаряется:
\[
\tau_{ij} \leftarrow (1-\rho)\tau_{ij},
\]
а затем на ребрах успешных путей добавляется новый феромон:
\[
\Delta\tau_{ij}=\frac{Q}{L}.
\]
Здесь $L$ --- длина найденного пути, $Q$ --- коэффициент количества феромона.

\section{Программная реализация}
Ниже приведен полный код программы на Python. В нем реализованы оба метода: генетический алгоритм для задачи коммивояжера и алгоритм муравьиной колонии для поиска кратчайшего пути.

\begin{lstlisting}[caption={Код программы для лабораторной работы №6}]
import random
import math

# ---------- Task 1: genetic algorithm for TSP ----------
DIST_TSP = [
    [0, 1, 11, 10, 5],
    [1, 0, 11, 6, 10],
    [11, 11, 0, 6, 1],
    [10, 6, 6, 0, 1],
    [5, 10, 1, 1, 0],
]


def route_length(genome):
    # City 1 is fixed as the start and finish.
    route = [0] + genome + [0]
    return sum(DIST_TSP[route[i]][route[i + 1]] for i in range(len(route) - 1))


def fitness(genome):
    return 1 / route_length(genome)


def order_crossover(parent1, parent2):
    n = len(parent1)
    left, right = sorted(random.sample(range(n), 2))
    child = [None] * n
    child[left:right + 1] = parent1[left:right + 1]

    pos = (right + 1) % n
    for gene in parent2:
        if gene not in child:
            child[pos] = gene
            pos = (pos + 1) % n
    return child


def mutate(genome, mutation_rate):
    if random.random() < mutation_rate:
        i, j = random.sample(range(len(genome)), 2)
        genome[i], genome[j] = genome[j], genome[i]


def select_parent(population):
    total = sum(fitness(g) for g in population)
    point = random.random() * total
    current = 0
    for genome in population:
        current += fitness(genome)
        if current >= point:
            return genome[:]
    return population[-1][:]


def genetic_algorithm_tsp(seed=11, population_size=4, mutation_rate=0.01, generations=200):
    random.seed(seed)
    genes = list(range(1, len(DIST_TSP)))

    population = []
    while len(population) < population_size:
        genome = genes[:]
        random.shuffle(genome)
        if genome not in population:
            population.append(genome)

    best = min(population, key=route_length)
    improvements = [(0, best[:], route_length(best))]

    for generation in range(1, generations + 1):
        new_population = [min(population, key=route_length)[:]]
        while len(new_population) < population_size:
            parent1 = select_parent(population)
            parent2 = select_parent(population)
            child = order_crossover(parent1, parent2)
            mutate(child, mutation_rate)
            new_population.append(child)

        population = new_population
        current_best = min(population, key=route_length)
        if route_length(current_best) < route_length(best):
            best = current_best[:]
            improvements.append((generation, best[:], route_length(best)))

    return best, route_length(best), improvements


# ---------- Task 2: ant colony algorithm for the shortest path ----------
NAMES = ["A", "B", "C", "D", "E", "F", "G"]
INF = math.inf
GRAPH = [
    [INF, 17, 5, INF, 8, INF, INF],
    [17, INF, INF, INF, INF, INF, INF],
    [5, INF, INF, 17, 17, INF, 45],
    [INF, INF, 17, INF, 9, 11, 22],
    [8, INF, 17, 9, INF, INF, INF],
    [INF, INF, INF, 11, INF, INF, 13],
    [INF, INF, 45, 22, INF, 13, INF],
]


def available_neighbors(vertex, visited):
    return [
        u for u, weight in enumerate(GRAPH[vertex])
        if not math.isinf(weight) and u not in visited
    ]


def build_ant_path(start, finish, pheromone, alpha, beta):
    path = [start]
    visited = {start}
    total_length = 0
    current = start

    while current != finish:
        candidates = available_neighbors(current, visited)
        if not candidates:
            return None, math.inf

        desirabilities = []
        for nxt in candidates:
            tau = pheromone[current][nxt] ** alpha
            eta = (1 / GRAPH[current][nxt]) ** beta
            desirabilities.append(tau * eta)

        border = random.random() * sum(desirabilities)
        cumulative = 0
        chosen = candidates[-1]
        for nxt, value in zip(candidates, desirabilities):
            cumulative += value
            if cumulative >= border:
                chosen = nxt
                break

        path.append(chosen)
        visited.add(chosen)
        total_length += GRAPH[current][chosen]
        current = chosen

    return path, total_length


def ant_colony_shortest_path(seed=11, iterations=80, ants_per_iteration=20,
                             alpha=1.0, beta=2.0, evaporation=0.3, q=100):
    random.seed(seed)
    start, finish = 0, 6
    n = len(GRAPH)
    pheromone = [
        [0.0 if math.isinf(GRAPH[i][j]) else 1.0 for j in range(n)]
        for i in range(n)
    ]

    best_path = None
    best_length = math.inf
    improvements = []

    for iteration in range(1, iterations + 1):
        successful_paths = []
        for _ in range(ants_per_iteration):
            path, length = build_ant_path(start, finish, pheromone, alpha, beta)
            if path is not None:
                successful_paths.append((path, length))
                if length < best_length:
                    best_path = path[:]
                    best_length = length
                    improvements.append((iteration, best_path[:], best_length))

        for i in range(n):
            for j in range(n):
                if not math.isinf(GRAPH[i][j]):
                    pheromone[i][j] *= (1 - evaporation)

        for path, length in successful_paths:
            delta = q / length
            for a, b in zip(path, path[1:]):
                pheromone[a][b] += delta
                pheromone[b][a] += delta

    return best_path, best_length, improvements


if __name__ == "__main__":
    best_genome, best_tsp_length, ga_improvements = genetic_algorithm_tsp()
    best_route = [0] + best_genome + [0]
    best_route_text = " -> ".join(str(city + 1) for city in best_route)

    print("TASK 1. GENETIC ALGORITHM FOR TSP")
    print(f"Best route: {best_route_text}")
    print(f"Route length: {best_tsp_length}")
    print("Improvements:")
    for generation, genome, length in ga_improvements:
        route = [0] + genome + [0]
        route_text = " -> ".join(str(city + 1) for city in route)
        print(f"  generation {generation}: {route_text}, length = {length}")

    best_path, best_path_length, aco_improvements = ant_colony_shortest_path()
    best_path_text = " -> ".join(NAMES[i] for i in best_path)

    print()
    print("TASK 2. ANT COLONY ALGORITHM FOR SHORTEST PATH")
    print(f"Best path: {best_path_text}")
    print(f"Path length: {best_path_length}")
    print("Improvements:")
    for iteration, path, length in aco_improvements:
        path_text = " -> ".join(NAMES[i] for i in path)
        print(f"  iteration {iteration}: {path_text}, length = {length}")

\end{lstlisting}

\section{Вывод программы}
После запуска программы был получен следующий результат:

\begin{lstlisting}[caption={Вывод программы},language={}]
TASK 1. GENETIC ALGORITHM FOR TSP
Best route: 1 -> 5 -> 3 -> 4 -> 2 -> 1
Route length: 19
Improvements:
  generation 0: 1 -> 2 -> 3 -> 4 -> 5 -> 1, length = 24
  generation 3: 1 -> 5 -> 3 -> 4 -> 2 -> 1, length = 19

TASK 2. ANT COLONY ALGORITHM FOR SHORTEST PATH
Best path: A -> E -> D -> G
Path length: 39
Improvements:
  iteration 1: A -> C -> E -> D -> F -> G, length = 55
  iteration 1: A -> E -> D -> F -> G, length = 41
  iteration 6: A -> E -> D -> G, length = 39
\end{lstlisting}

\section{Анализ результатов}
\subsection{Результат для задачи коммивояжера}
Генетический алгоритм нашел маршрут:
\[
1 \rightarrow 5 \rightarrow 3 \rightarrow 4 \rightarrow 2 \rightarrow 1.
\]
Длина маршрута равна:
\[
5 + 1 + 6 + 6 + 1 = 19.
\]
Так как матрица расстояний симметрична, обратный маршрут
\[
1 \rightarrow 2 \rightarrow 4 \rightarrow 3 \rightarrow 5 \rightarrow 1
\]
имеет такую же длину. Полученное значение целевой функции равно 19, то есть именно эту сумму расстояний алгоритм минимизировал.

\subsection{Результат для алгоритма муравьиной колонии}
Алгоритм муравьиной колонии нашел путь:
\[
A \rightarrow E \rightarrow D \rightarrow G.
\]
Его длина равна:
\[
8 + 9 + 22 = 39.
\]
В процессе работы сначала были найдены более длинные пути: $A \rightarrow C \rightarrow E \rightarrow D \rightarrow F \rightarrow G$ длиной 55 и $A \rightarrow E \rightarrow D \rightarrow F \rightarrow G$ длиной 41. После накопления феромона на более выгодных ребрах был найден путь длиной 39.

\section{Вывод}
В ходе лабораторной работы были решены две задачи оптимизации для варианта 11. Для задачи коммивояжера был применен генетический алгоритм с размером популяции 4 и вероятностью мутации 0,01. Лучший найденный маршрут имеет длину 19.

Для задачи поиска кратчайшего пути был применен алгоритм муравьиной колонии. Лучший найденный путь из A в G имеет вид $A \rightarrow E \rightarrow D \rightarrow G$, его суммарный вес равен 39.

Таким образом, оба эвристических метода позволили получить оптимальные или близкие к оптимальным решения за небольшое число итераций. Генетический алгоритм оказался удобным для работы с перестановками городов, а алгоритм муравьиной колонии --- для выбора последовательности ребер в графе с учетом накопленного феромона и расстояний.

\end{document}
